\documentclass[11pt]{article}
\usepackage[a4paper,margin=2.2cm]{geometry}
\usepackage{xcolor}
\usepackage{fontspec}
\usepackage{xeCJK}
\usepackage{enumitem}
\setmainfont{Times New Roman}
\setCJKmainfont{SimSun}
\setlist[itemize]{leftmargin=1.4em,itemsep=0.35em,topsep=0.35em}
\setlength{\parindent}{0pt}
\setlength{\parskip}{0.55em}
\pagestyle{plain}
\begin{document}
{\color[HTML]{1F5FD2}\LARGE\bfseries 融合课堂教学策略与调整\par}
{\color[HTML]{667085}\small 星轨原创教育资源：用于家庭与学校共同制定支持计划。\par}
\vspace{0.8em}
{\color[HTML]{111827}\large\bfseries 通用学习设计\par}
\begin{itemize}
\item 融合课堂应预先设计多种接收信息、表达学习和保持参与的方式，而不是等儿童失败后再临时补救。
\item 同一学习目标可以通过阅读、图片、实物、讨论、操作或短视频进入，降低单一路径带来的障碍。
\end{itemize}
{\color[HTML]{111827}\large\bfseries 差异化教学\par}
\begin{itemize}
\item 内容差异化可以使用不同阅读难度、图片、实物或预教词汇。
\item 过程差异化可以安排独立练习、小组合作、同伴辅导、学习站和教师指导练习。
\item 成果差异化可允许口头报告、绘画、模型、短写作、录音或演示，只要能体现核心目标。
\end{itemize}
{\color[HTML]{111827}\large\bfseries 同伴支持\par}
\begin{itemize}
\item 同伴支持应是互惠的。教师要轮换伙伴、明确角色，并教全班如何尊重地协作。
\item 不要让同伴长期承担成人职责，成人仍需负责观察、指导和及时调整支持。
\end{itemize}
\vfill
{\color[HTML]{667085}\footnotesize 本材料为应用内原创教育内容，不能替代医疗、法律或正式评估建议。\par}
\end{document}
